\section{Proofs for Normal Domains of Attraction}
\label{sec:annex-p15-normal-domains-of-attraction}

\begin{annexitem}{annex:p15.nda-gk-truncated-mean}{thm:p15-nda-gk-truncated-mean}{Constructing the source weight}
\begin{proof}
Choose strictly increasing \(r_k\geq k\), \(k\geq1\), so large that
\(\mathbb P(|X|>r_k)\leq2^{-k}/(2k+1)\).
This is possible because \(X\) is finite.
Define the finite Borel function
\[
 w(x)=1+\sum_{k=1}^\infty\mathbf1_{\{|x|>r_k\}}.
 \]
It is at least one and tends to infinity as \(|x|\to\infty\).
Since the indicator sets decrease with \(k\), its square has
the exact nonnegative expansion
\[
 w(x)^2=1+\sum_{k=1}^\infty(2k+1)\mathbf1_{\{|x|>r_k\}}.
 \]
Monotone convergence therefore gives
\(\mathbb Ew(X)^2\leq1+\sum_{k\geq1}2^{-k}=2\).
For \(U(C)=\mathbb E(X^2;|X|\leq C)\), Cauchy--Schwarz yields
\[
 |m(C)|^2
 =\left|\mathbb E\left(w(X)\frac{X}{w(X)};|X|\leq C\right)\right|^2
 \leq2\,\mathbb E\left(\frac{X^2}{w(X)^2};|X|\leq C\right).
 \]
Fix an integer \(K\geq1\). Outside a sufficiently large bounded
interval, \(w\geq K\), whereas the integral of \(X^2/w^2\)
over that interval is a finite constant \(D_K\).
Thus the final expectation is at most \(D_K+U(C)/K^2\).
Since \(U(C)\to\infty\),
\(\limsup_Cm(C)^2/U(C)\leq2/K^2\).
Let \(K\to\infty\). The variance identity
\(\operatorname{Var}(X\mathbf1_{\{|X|\leq C\}})
=U(C)-m(C)^2\) proves the second assertion.
This supplies the construction and the limiting step behind the
source's weighted Cauchy--Schwarz argument.
\end{proof}
\end{annexitem}

\begin{annexitem}{annex:p15.nda-gk-cutoff-interface}{thm:p15-nda-gk-cutoff-interface}{The whole i.i.d. cutoff premise}
\begin{proof}
For sufficiency, put
\(V_n=\operatorname{Var}(X\mathbf1_{\{|X|\leq C_n\}})\)
and \(B_n^2=nV_n\). The assumption \(nV_n/C_n^2\to\infty\)
implies \(B_n>0\) eventually and \(C_n/B_n\to0\).
The independent centered row
\[
 \frac{X_j\mathbf1_{\{|X_j|\leq C_n\}}-m(C_n)}{B_n},
                       \qquad1\leq j\leq n,
 \]
has variance sum one and absolute bound \(2C_n/B_n\to0\),
because \(|m(C_n)|\leq C_n\).
Its Lindeberg tails vanish eventually at every positive cutoff,
so it converges to \(\chi\). With probability at least
\(1-nq_X(C_n)\), no summand is removed. The same normal limit
therefore holds for \((S_n-nm(C_n))/B_n\).
Also \(B_n\to\infty\), and the i.i.d. individual summands
\(X_j/B_n\), \(1\leq j\leq n\), are infinitesimal, meaning
\(\max_{j\leq n}\mathbb P(|X_j/B_n|>\varepsilon)\to0\)
for every \(\varepsilon>0\).

For necessity, suppose normal attraction holds.
The independent-difference argument in
Theorem~\ref{thm:p15-nda-characterization}, the symmetric-kernel
proof, and Theorem~\ref{thm:p15-nda-desymmetrization} give the
tail criterion and slow variation of \(U=H_X\).
They did not use this cutoff interface.
The moment theorem gives the mean \(\mu\), and because
\(\mathbb EX^2=\infty\), the centered truncated moment
\(H_{X-\mu}(x)\sim U(x)\).
Let \(b_n\) be the continuous scales already constructed from this
centered law; their construction, as opposed to the sufficient CLT,
implies for each fixed integer \(k\geq1\),
\[
 nq_X(b_n/k)\to0,\qquad
 \frac{nU(b_n/k)}{b_n^2}\to1,\qquad b_n/k\to\infty.
 \]
Here the first limit follows by multiplying the tail ratio at
\(b_n/k\) by \(k^2nU(b_n/k)/b_n^2\).
Choose increasing integers \(N_k\) so that, for all \(n\geq N_k\),
\[
 nq_X(b_n/k)\leq1/k,\qquad
 nU(b_n/k)/b_n^2\geq1/2,\qquad b_n/k\geq k,
 \]
and, by the preceding weighted-mean lemma,
\(\operatorname{Var}(X\mathbf1_{\{|X|\leq b_n/k\}})
\geq U(b_n/k)/2\).
Set \(C_n=b_n/k\) for \(N_k\leq n<N_{k+1}\);
choose any positive cutoffs before \(N_1\).
Then \(C_n\to\infty\), \(nq_X(C_n)\to0\), and on that block
\[
 \frac n{C_n^2}\operatorname{Var}(X\mathbf1_{\{|X|\leq C_n\}})
 \geq k^2/4\longrightarrow\infty.
 \]
The source weighted-mean lemma also proves equivalence with
\(nU(C_n)/C_n^2\to\infty\).
This proves the full required i.i.d. interface, with its
necessity supplied by a separately identified argument.
\end{proof}
\end{annexitem}

\begin{annexitem}{annex:p15.nda-gk-cutoff-route}{thm:p15-nda-gk-cutoff-route}{The source's two cutoff implications}
\begin{proof}
We retain the source's scaled truncated moment
\(\mathcal H(x)=U(x)/x^2\), distinguishing it from the
unscaled moment \(H_X\) used earlier.
The exact identity
\[
 \mathcal H(x)+q(x)
       =\mathbb E\min(X^2/x^2,1)
 \]
makes this sum nonincreasing in \(x>0\).
Suppose \(C_n\to\infty\),
\(nq(C_n)\to0\), and \(n\mathcal H(C_n)\to\infty\).
For sufficiently large \(x\), set
\(n+1=\min\{j\geq1:C_j>x\}\), which is eventually at least two.
Then \(C_n\leq x<C_{n+1}\). Such indices tend to
infinity as \(x\to\infty\), even if the sequence is not monotone.
Monotonicity of \(q\) and of \(\mathcal H+q\) gives
\[
 q(x)\leq q(C_n),\qquad
 \mathcal H(x)\geq\mathcal H(C_{n+1})-q(C_n).
 \]
The latter right side is positive for large \(n\), since
\[
 n\mathcal H(C_{n+1})-nq(C_n)
 =\frac n{n+1}\{(n+1)\mathcal H(C_{n+1})\}-nq(C_n)
 \longrightarrow\infty.
 \]
Consequently
\[
 0\leq\frac{q(x)}{\mathcal H(x)}
 \leq\frac{nq(C_n)}
       {\frac n{n+1}\{(n+1)\mathcal H(C_{n+1})\}-nq(C_n)}
 \longrightarrow0.
 \]
This is the source's cutoff sandwich, with its positive
denominator and first-crossing index fully specified.

For the converse fix \(\delta>0\).
The tail \(q(x)=\mathbb P(|X|>x)\) is nonincreasing and
right-continuous, and tends to zero. Therefore
\[
 C_n(\delta)=\inf\{x\geq0:nq(x)\leq\delta\}
 \]
is finite, and its right-continuity ensures
\(nq(C_n(\delta))\leq\delta\).
An infinite second moment forces unbounded support, so for
every fixed \(M>0\), \(q(M)>0\); hence \(C_n(\delta)>M\)
eventually and \(C_n(\delta)\to\infty\).
At half the cutoff, for all sufficiently large \(n\),
\[
 nq(C_n(\delta)/2)>\delta.
 \]
Since \(q(x)/\mathcal H(x)\to0\), this gives
\(n\mathcal H(C_n(\delta)/2)\to\infty\).
Also monotonicity of \(U\) gives
\(\mathcal H(C_n(\delta))\geq
\mathcal H(C_n(\delta)/2)/4\).
Thus \(n\mathcal H(C_n(\delta))\to\infty\) for every fixed
\(\delta>0\), exactly the intermediate source conclusion.

Choose increasing integers \(M_k\) so that for all \(n\geq M_k\),
\[
 C_n(1/k)\geq k,\qquad n\mathcal H(C_n(1/k))\geq k.
 \]
Set \(\delta_n=1/k\) on \(M_k\leq n<M_{k+1}\), and set
\(C_n=C_n(\delta_n)\) there, with arbitrary positive initial
values. Then \(C_n\to\infty\),
\(nq(C_n)\leq1/k\to0\), and
\(n\mathcal H(C_n)\geq k\to\infty\).
The weighted-mean lemma replaces \(U(C_n)\) by the truncated
variance, and the proved cutoff interface supplies the source's
normal limit. This preserves its tail-quantile diagonal construction,
which is different from the earlier continuous winsorized inverse.

For completeness, if \(X\) is nondegenerate with finite variance
\(\sigma^2>0\), dominated convergence gives
\[
 x^2q_X(x)\leq\mathbb E(X^2;|X|>x)\to0,\qquad
 H_X(x)\to\mathbb EX^2>0.
 \]
After centering by \(\mu\), the entries
\((X_j-\mu)/(\sigma\sqrt n)\) have variance sum one and
Lindeberg tail
\(\sigma^{-2}\mathbb E((X-\mu)^2;
|X-\mu|>\varepsilon\sigma\sqrt n)\to0\).
This is the sufficient case of Section 21, Theorem 3 used by
the source, supplied in full by the already proved general
Lindeberg theorem. It completes the source's finite-variance
branch and hence both branches of its Section 35 argument.
\end{proof}
\end{annexitem}

\begin{annexitem}{annex:p15.nda-tail-equivalence}{thm:p15-nda-tail-equivalence}{Dyadic shells and the exact cutoff conventions}
\begin{proof}
Put \(R=|Z|\), \(H=H_Z\), and \(q=q_Z\).
If \(x^2q(x)=o(H(x))\), then for each fixed \(c>1\),
\[
 0\leq H(cx)-H(x)
 =\mathbb E(R^2;x<R\leq cx)
 \leq c^2x^2q(x)=o(H(x)).
 \]
This gives \(H(cx)/H(x)\to1\) for \(c>1\).
Replacing \(x\) by \(cx\) and taking reciprocals gives the
same assertion for \(0<c<1\). The case \(c=1\) is immediate.

Conversely suppose \(H\) is slowly varying. We first establish
the only uniform bound on slow variation needed in this chapter.
Fix \(0<\delta<2\). There is \(x_0>0\), with \(H(x_0)>0\),
such that
\[
 H(2u)\leq2^\delta H(u)\quad(u\geq x_0).
 \]
Iterating gives \(H(2^ku)\leq2^{k\delta}H(u)\).
For \(v\geq u\geq x_0\), choose the integer \(k\geq0\)
with \(2^ku\leq v<2^{k+1}u\). Monotonicity therefore gives
\[
 H(v)\leq2^{(k+1)\delta}H(u)
                    \leq2^\delta(v/u)^\delta H(u).                 \tag{*}
 \]
In particular \(H(v)\leq C_\delta v^\delta\) for large \(v\),
where \(C_\delta=2^\delta H(x_0)/x_0^\delta\).

Because \(R\) is finite, the disjoint shells
\((2^jx,2^{j+1}x]\), \(j\geq0\), exhaust \(\{R>x\}\).
On the \(j\)-th shell \(R^2\geq 2^{2j}x^2\), so
\[
 \frac{x^2q(x)}{H(x)}
 \leq\sum_{j=0}^\infty
 \frac{H(2^{j+1}x)-H(2^jx)}{2^{2j}H(x)}.
 \]
For each fixed \(j\), the summand tends to zero by slow variation.
For \(x\geq x_0\), it is bounded above by
\(2^{(j+1)\delta-2j}\), a summable geometric sequence.
To justify passage through the infinite sum, choose a finite
number of shells making this geometric majorant's remaining
tail arbitrarily small, then take the limit in the finite initial
sum. This proves \(x^2q(x)/H(x)\to0\).

The exact decomposition
\(W_Z(x)=H(x)+x^2q(x)\) follows by partitioning according
to \(R\leq x\) and \(R>x\). It proves \(W_Z\sim H\),
and thus \(W_Z\) is slowly varying.
The mass at a cutoff obeys
\[
 x^2\mathbb P(R=x)=H(x)-H(x-)
                      \leq H(x)-H(x/2)=o(H(x)).
 \]
Hence \(H(x-)/H(x)\to1\), and adding this atom to \(q(x)\)
proves the closed-tail assertion. All formulas have therefore
retained their strict or closed event conventions.
\end{proof}
\end{annexitem}

\begin{annexitem}{annex:p15.nda-moments-and-shifts}{thm:p15-nda-moments-and-shifts}{Tail moments and fixed changes of origin}
\begin{proof}
Continue with \(R=|Z|\) and \(H=H_Z\).
For \(x\) sufficiently large define
\[
 e(x)=\sup_{u\geq x}\frac{u^2q_Z(u)}{H(u)}.
 \]
The ratio tends to zero, so \(e(x)\) is finite eventually and
decreases to zero. For \(0<p<2\), choose
\(0<\delta<2-p\) and use bound \((*)\) from the preceding proof.
For \(u\geq x\) sufficiently large,
\[
 q_Z(u)\leq e(x)\frac{H(u)}{u^2}
 \leq2^\delta e(x)H(x)x^{-\delta}u^{\delta-2}.
 \]
The pointwise layer identity for \(R^p\) on \(\{R>x\}\)
and Tonelli give
\[
 \mathbb E(R^p;R>x)
 =x^p q_Z(x)+p\int_x^\infty u^{p-1}q_Z(u)\,du
 \leq e(x)H(x)x^{p-2}
       \left(1+\frac{p2^\delta}{2-p-\delta}\right).
 \]
The integral converges because \(p+\delta<2\).
The bounded portion \(R\leq x\) has finite \(p\)-moment, so this
also proves full moment finiteness and the claimed little-o bound.
In particular \(\mathbb E|Z|<\infty\).

Fix \(c\in\mathbb R\), and set \(a=|c|\).
If \(\mathbb EZ^2<\infty\), then \(\mathbb E(Z-c)^2<\infty\).
Nondegeneracy makes both second moments positive.
Their truncated second moments tend to these positive constants,
and are therefore slowly varying. The preceding equivalence
gives both tail criteria.

Suppose instead that \(\mathbb EZ^2=\infty\), so \(H(x)\to\infty\).
For \(x>a\),
\[
 \{|Z|\leq x-a\}\subseteq\{|Z-c|\leq x\}
                     \subseteq\{|Z|\leq x+a\}.
 \]
The difference of the two squared functions is integrable:
\[
 \mathbb E\left|(Z-c)^2-Z^2\right|
 \leq2a\,\mathbb E|Z|+a^2=:C<\infty.
 \]
On each bounded cutoff event both squared integrals are finite,
and their difference is therefore bounded in absolute value by \(C\).
Consequently
\[
 H(x-a)-C\leq H_{Z-c}(x)\leq H(x+a)+C.
 \]
Monotonicity and slow variation show
\(H(x\pm a)/H(x)\to1\): for any fixed \(0<r<1<s\),
eventually \(rx\leq x-a\leq x+a\leq sx\), and the outer
ratios tend to one. Divide the preceding sandwich by \(H(x)\).
It follows that \(H_{Z-c}(x)\sim H(x)\), hence is slowly varying.
The tail equivalence completes the shift claim. In particular,
centering by the now-existing mean gives the stated function \(V\).
\end{proof}
\end{annexitem}

\begin{annexitem}{annex:p15.nda-normalizer-construction}{thm:p15-nda-normalizer-construction}{Inverting a continuous decreasing function}
\begin{proof}
Write \(W=W_Y\), \(V=H_Y\), and
\[
 g(x)=W(x)/x^2
     =\mathbb E\min(Y^2/x^2,1)\quad(x>0).
 \]
For each finite \(Y\), the integrand is continuous in \(x>0\),
bounded by one, and tends to zero as \(x\to\infty\).
Bounded convergence proves continuity of \(g\) and \(g(x)\to0\).
It is nonincreasing. For \(x_0\leq x_1<x_2\), the
integrand decreases strictly on
\(\{0<|Y|\leq x_0\}\), an event of positive probability,
so \(g(x_1)>g(x_2)\).
Thus for every sufficiently large \(n\), with
\(0<1/n<g(x_0)\), the intermediate value theorem gives one
and only one root \(b_n>x_0\). The roots increase and tend
to infinity, since \(g\) is continuous and positive at every
fixed finite \(x\geq x_0\).

By the tail equivalence, \(W(x)\sim V(x)\).
The defining equation therefore gives
\[
 nV(b_n)/b_n^2=V(b_n)/W(b_n)\to1.
 \]
For each fixed \(\varepsilon>0\), slow variation gives
\(V(\varepsilon b_n)/V(b_n)\to1\).
Multiplying the tail ratio at \(\varepsilon b_n\) by
\(nV(\varepsilon b_n)/(\varepsilon^2b_n^2)\)
shows \(nq_Y(\varepsilon b_n)\to0\).
The preceding tail-moment bound at \(p=1\) gives
\[
 \frac{n}{b_n}\mathbb E(|Y|;|Y|>b_n)
 =\frac{nV(b_n)}{b_n^2}\,
      \frac{b_n\mathbb E(|Y|;|Y|>b_n)}{V(b_n)}\to0.
 \]

Finally \(g(rx)/g(x)=r^{-2}W(rx)/W(x)\to r^{-2}\)
for each fixed \(r>0\). Fix \(c>0\), and choose
\(0<r<\sqrt c<s\).
The defining roots obey, for all sufficiently large \(n\),
\[
 g(rb_n)>1/\lfloor cn\rfloor>g(sb_n),
 \]
because \(ng(rb_n)\to r^{-2}>c^{-1}\) and
\(ng(sb_n)\to s^{-2}<c^{-1}\).
Strict decrease of \(g\) gives
\(r<b_{\lfloor cn\rfloor}/b_n<s\).
Let \(r\uparrow\sqrt c\) and \(s\downarrow\sqrt c\).
This proves the asserted scale ratio and completes the construction.
\end{proof}
\end{annexitem}

\begin{annexitem}{annex:p15.nda-sufficiency}{thm:p15-nda-sufficiency}{The complete Lindeberg and centering calculation}
\begin{proof}
The tail and shift results give the finite mean \(\mu=\mathbb EX\)
and slow variation of \(V=H_Y\), \(Y=X-\mu\).
Let \(Y_j=X_j-\mu\), and define
\[
 T_{n,j}=Y_j\mathbf1_{\{|Y_j|\leq b_n\}},\quad
 m_n=\mathbb ET_{n,1},\quad
 U_{n,j}=(T_{n,j}-m_n)/b_n\quad(1\leq j\leq n).
 \]
These rows are independent and centered.
Since \(\mathbb EY=0\),
\[
 |m_n|=|\mathbb E(Y;|Y|>b_n)|
 \leq\mathbb E(|Y|;|Y|>b_n),\qquad
 n|m_n|/b_n\to0.
 \]
In particular \(nm_n^2/b_n^2=(n|m_n|/b_n)^2/n\to0\).
The row variance sum is therefore
\[
 \sum_{j=1}^n\mathbb EU_{n,j}^2
   =\frac n{b_n^2}\{V(b_n)-m_n^2\}\to1.
 \]

Fix \(\varepsilon>0\) and put
\(r=\min(\varepsilon/2,1/2)>0\).
Eventually \(|m_n|\leq\varepsilon b_n/2\).
On \(|T_{n,j}-m_n|>\varepsilon b_n\), the truncated variable
must satisfy \(|T_{n,j}|>\varepsilon b_n/2\), and hence
\(rb_n<|Y_j|\leq b_n\). Using
\((a-b)^2\leq2a^2+2b^2\),
\[
 \sum_{j=1}^n\mathbb E(U_{n,j}^2;|U_{n,j}|>\varepsilon)
 \leq \frac{2n}{b_n^2}\{V(b_n)-V(rb_n)\}
                              +\frac{2nm_n^2}{b_n^2}\to0.
 \]
The last limit follows from slow variation and
\(nV(b_n)/b_n^2\to1\). Thus the full Lindeberg--Feller theorem
(Theorem~\ref{thm:p15-clt-lindeberg-feller}) gives
\(\sum_{j=1}^nU_{n,j}\Rightarrow\chi\).

The union bound and the constructed scale give
\[
 \mathbb P(\exists j\leq n:Y_j\ne T_{n,j})
 \leq nq_Y(b_n)\to0.
 \]
Except on that event, the difference between
\(\sum_jY_j/b_n\) and \(\sum_jU_{n,j}\) is the deterministic
quantity \(nm_n/b_n\to0\).
Therefore the difference tends to zero in probability, and the
proved perturbation lemma transfers the normal limit.
This establishes the sufficient condition with exact centering
\(n\mu\), for finite or infinite variance.
\end{proof}
\end{annexitem}

\begin{annexitem}{annex:p15.nda-normal-types}{thm:p15-nda-normal-types}{Two fixed quantiles determine location and scale}
\begin{proof}
Let \(F_n\) be the CDF of \(Z_n\) and
\(Q_n(p)=\inf\{x:F_n(x)\geq p\}\) for \(0<p<1\).
For a centered nondegenerate normal limit with CDF \(G\),
let \(q(p)\) be its unique \(p\)-quantile.
Its density is strictly positive, so for every \(\varepsilon>0\),
\[
 G(q(p)-\varepsilon)<p<G(q(p)+\varepsilon).
 \]
Weak convergence at these continuity points gives the same strict
inequalities with \(F_n\) eventually. The elementary quantile
inequalities then place \(Q_n(p)\) in
\([q(p)-\varepsilon,q(p)+\varepsilon]\).
Hence \(Q_n(p)\to q(p)\).

A positive affine transformation sends the \(p\)-quantile exactly
to \(c_nQ_n(p)+d_n\), by its CDF
\(F_n((x-d_n)/c_n)\) and the definition of the infimum.
The assumed second weak limit therefore gives
\(c_nQ_n(p)+d_n\to q(p)\).
Take \(p=1/4\) and \(p=3/4\), and subtract.
The limiting difference \(q(3/4)-q(1/4)>0\), so \(c_n\to1\).
Substitution at either quantile gives \(d_n\to0\).
\end{proof}
\end{annexitem}

\begin{annexitem}{annex:p15.nda-symmetric-necessity}{thm:p15-nda-symmetric-necessity}{The nonnegative kernel and passage to all scales}
\begin{proof}
Write \(T_n=\sum_{j=1}^nZ_j\),
\(\psi(t)=\mathbb Ee^{itZ_1}=\mathbb E\cos(tZ_1)\),
and \(H=H_Z\), \(q=q_Z\).
We first show \(b_n\to\infty\).
Nondegeneracy and symmetry imply
\(\mathbb P(0<|Z_1|\leq M)>0\) for some finite \(M>0\).
For \(0<|t|<2\pi/M\), the integrand \(1-\cos(tZ_1)\)
is strictly positive on this event, so \(\psi(t)<1\).
Continuity at zero also makes \(\psi(t)>0\) on some smaller
neighborhood of zero. Thus \(0<\psi(t)<1\) for every nonzero
\(t\) in a fixed sufficiently small neighborhood.

If \(b_n\) had a bounded subsequence, tightness of \(T_n/b_n\)
would make the unnormalized \(T_n\) tight on that subsequence:
for an upper bound \(B\) on \(b_n\),
\(\mathbb P(|T_n|>BK)\leq\mathbb P(|T_n/b_n|>K)\).
Helly selection and tightness would give a further proper weak
limit. But on the punctured neighborhood just found its
characteristic function would be
\(\lim\psi(t)^n=0\), while its value at zero must be one.
This contradicts continuity of characteristic functions.
Hence \(b_n\to\infty\).

Fix \(A>0\). The transforms of \(T_n/b_n\) converge uniformly
on \([-2A,2A]\) to \(e^{-\tau^2t^2/2}\), by
Theorem~\ref{thm:p15-criteria-tight-transforms}.
Also \(\psi(t/b_n)\to1\) uniformly there by continuity at zero,
so eventually it is positive throughout that interval.
The limiting Gaussian transform is bounded away from zero there.
Taking the real logarithm therefore gives uniformly
\[
 -n\log\psi(t/b_n)\longrightarrow\tau^2t^2/2.
 \]
For \(0\leq u\leq1/2\),
\(0\leq-\log(1-u)-u\leq u^2\), as proved by integrating
\(u/(1-u)\). The preceding uniform logarithm bound makes
\(n(1-\psi(t/b_n))\) bounded uniformly; multiplying by the
uniformly vanishing \(1-\psi(t/b_n)\) controls the remainder.
Consequently, uniformly on this interval,
\[
 n\{1-\psi(t/b_n)\}\longrightarrow\tau^2t^2/2.             \tag{1}
 \]

The identity
\[
 4(1-\cos u)-(1-\cos2u)=2(1-\cos u)^2
 \]
and (1) imply
\[
 n\mathbb E(1-\cos(tZ_1/b_n))^2\longrightarrow0
                                  \quad\text{uniformly for }|t|\leq A.
 \]
Integrate over \(0\leq t\leq A\); Tonelli applies to this
nonnegative integrand. Define
\[
 K_A(x)=\int_0^A(1-\cos(tx))^2\,dt.
 \]
This function is continuous, is positive for every \(x\ne0\),
and, for \(x\ne0\), equals
\[
 K_A(x)=\frac{3A}{2}-\frac{2\sin(Ax)}x+\frac{\sin(2Ax)}{4x}.
 \]
The formula follows by expanding the square and integrating the
cosines. It shows \(K_A(x)\to3A/2\) as \(|x|\to\infty\).
Its positivity follows because an integral of a nonnegative
continuous function can be zero only when that function vanishes
identically, which would force \(\cos(tx)=1\) for all \(t\)
in an interval and hence \(x=0\).
For each \(r>0\), it follows that
\(\inf_{|x|\geq r}K_A(x)>0\): use the positive limiting value
outside a large compact set and the positive minimum on the
remaining compact annulus. We have proved
\(n\mathbb EK_A(Z_1/b_n)\to0\); hence
\[
 n q(rb_n)\to0\quad(r>0).                                \tag{2}
 \]
This excludes jumps even at frequencies where an individual
cosine alone would fail to see them.

Fix \(r>0\). By (2), the contribution of \(|Z_1|>rb_n\)
to \(n\mathbb E[1-\cos(tZ_1/b_n)]\) is at most
\(2nq(rb_n)\to0\). For \(t\ne0\) so small that
\(t^2r^2<12\), Taylor's integral or alternating-series bounds give
\[
 \frac{t^2x^2}{2}\left(1-\frac{t^2r^2}{12}\right)
 \leq1-\cos(tx)\leq\frac{t^2x^2}{2}\quad(|x|\leq r).
 \]
For clarity, for \(u\geq0\) Taylor's exact integral remainder gives
\[
 \cos u=1-u^2/2+
              \frac16\int_0^u(u-s)^3\cos s\,ds
          \leq1-u^2/2+u^4/24.
 \]
Evenness proves the same bound for negative \(u\). The upper bound
\(1-\cos u\leq u^2/2\) follows by integrating
\(|\sin s|\leq|s|\). Thus neither bound needs an unproved
uniform approximation.
Let \(M_n(r)=nH(rb_n)/b_n^2\).
Combining (1), the negligible outside contribution, and these
bounds yields
\[
 \tau^2\leq\liminf_n M_n(r)
 \leq\limsup_n M_n(r)
 \leq\frac{\tau^2}{1-t^2r^2/12}.
 \]
Let \(t\to0\), keeping \(r\) fixed. Thus
\[
 nH(rb_n)/b_n^2\to\tau^2\quad(r>0).                       \tag{3}
 \]

Because \(b_n\to\infty\), the single variable \(Z_{n+1}/b_n\)
tends to zero in probability. Perturbation gives
\(T_{n+1}/b_n\Rightarrow N(0,\tau^2)\), while by assumption
\(T_{n+1}/b_{n+1}\Rightarrow N(0,\tau^2)\).
The affine-normalization lemma applied to their deterministic
scale ratio proves \(b_{n+1}/b_n\to1\).

It remains to deduce a statement for every large real \(x\),
not only for the sequence of normalizers.
Fix an initial index \(N\). For sufficiently large \(x\), let
\[
 m(x)=\min\{n\geq N:b_n\geq x\}.
 \]
The minimum exists because \(b_n\to\infty\), and \(m(x)\to\infty\).
Eventually \(m(x)>N\), and the first-crossing property gives
\(b_{m(x)-1}<x\leq b_{m(x)}\). Thus
\[
 1\leq b_{m(x)}/x
       <b_{m(x)}/b_{m(x)-1}\longrightarrow1.
 \]
For every fixed \(c>0\), monotonicity of \(H\) and (3),
applied at two fixed positive multiples bracketing
\(cx/b_{m(x)}\to c\), imply
\[
 \frac{m(x)}{b_{m(x)}^2}H(cx)\longrightarrow\tau^2.
 \]
The same holds with \(c=1\). Dividing proves
\(H(cx)/H(x)\to1\), the full slow-variation conclusion.
The tail-equivalence theorem now gives the ratio condition.
This passage does not require monotonicity of the original \(b_n\).
\end{proof}
\end{annexitem}

\begin{annexitem}{annex:p15.nda-desymmetrization}{thm:p15-nda-desymmetrization}{Two tail bounds and a truncated-moment comparison}
\begin{proof}
Write \(q_X=q\), and use subscripts \(Z\) for the independent
difference \(Z=X-X'\). Since \(X'\) is finite,
\(\mathbb P(|X'|\leq x/2)\geq1/2\) for all sufficiently large \(x\).
The independent event
\(\{|X|>x,\ |X'|\leq x/2\}\) implies \(|Z|>x/2\).
Therefore
\[
 q(x)\leq2q_Z(x/2).
 \]
Slow variation of \(H_Z\) and the tail equivalence imply
\[
 x^2q(x)=o(H_Z(x/2))=o(H_Z(x)).                           \tag{4}
 \]

On \(\{|Z|\leq x,\ |X|,|X'|\leq2x\}\),
\(Z^2\leq2X^2+2(X')^2\); its expectation is at most
\(4H_X(2x)\). On the rest of \(\{|Z|\leq x\}\),
at least one of \(|X|,|X'|\) exceeds \(2x\); the other must
exceed \(x\), by the triangle inequality.
There \(Z^2\leq x^2\), so independence and the union bound give
\[
 H_Z(x)\leq4H_X(2x)+2x^2q(2x)q(x).
 \]
By (4) and \(q(2x)\leq1\), the final term is \(o(H_Z(x))\).
Thus for all sufficiently large \(x\),
\(H_X(2x)\geq H_Z(x)/8\), and after replacing \(x\) by \(x/2\),
\[
 H_X(x)\geq H_Z(x/2)/8.
 \]
Combine this with (4) to obtain
\(x^2q_X(x)/H_X(x)\to0\).
Every expectation used before this conclusion was truncated and finite.
\end{proof}
\end{annexitem}

\begin{annexitem}{annex:p15.nda-characterization}{thm:p15-nda-characterization}{Both directions and all admissible centers and scales}
\begin{proof}
Equivalence of the tail criterion and slow variation was proved
above, and the sufficient condition gives the normal limit with
center \(n\mu\) and constructed scale \(b_n\).
It remains to prove necessity from an arbitrary presumed normal limit.

Suppose \((S_n-a_n)/d_n\Rightarrow\chi\), where \(d_n>0\).
On an independent product space construct a copy \(X'_j\)
of the whole sequence, and let \(S'_n=\sum_{j=1}^nX'_j\).
The two normalized centered sums are independent within each row.
Their difference has characteristic function the squared modulus
of the original transform, tending to \(e^{-t^2}\).
The characteristic-function continuity theorem gives
\[
 \frac{S_n-S'_n}{d_n}\Rightarrow N(0,2).
 \]
The differences \(Z_j=X_j-X'_j\) are i.i.d. and symmetric.
They are nondegenerate: if an independent difference were a
constant, symmetry would make it zero. For any Borel set \(A\),
\(X=X'\) almost surely and independence would then give
\(\mathbb P(X\in A)\{1-\mathbb P(X\in A)\}=0\).
In particular the CDF takes only zero and one values, and
right continuity and its endpoint limits force a single point mass.
That contradicts the assumed nondegeneracy of \(X\).
The symmetric-necessity theorem therefore makes \(H_Z\)
slowly varying. Desymmetrization proves the required tail ratio
for \(X\), without an assumed first or second moment.

The moment and shift theorem now supplies \(\mu\), and the
constructed normalization yields
\[
 U_n=(S_n-n\mu)/b_n\Rightarrow\chi.
 \]
For arbitrary deterministic \(a_n,d_n\), \(d_n>0\),
\[
 (S_n-a_n)/d_n=(b_n/d_n)U_n+(n\mu-a_n)/d_n.
 \]
If this also converges to \(\chi\), the affine-normalization
lemma gives \(b_n/d_n\to1\) and
\((n\mu-a_n)/d_n\to0\).
These are equivalent to the two displayed conditions in the theorem.
Conversely those conditions and the perturbation lemma give
the standard normal limit in this identity.
This proves necessity, sufficiency, construction, and uniqueness
of normalization up to asymptotically negligible changes.
\end{proof}
\end{annexitem}

\begin{annexitem}{annex:p15.nda-examples}{thm:p15-nda-examples}{The four boundary calculations}
\begin{proof}
For finite positive variance, the law of \(X\) is nondegenerate,
since a law concentrated at one point has variance zero.
The mean \(\mu=\mathbb EX\) exists, and \(Y=X-\mu\) is centered
and nondegenerate. By monotone convergence
\(H_Y(x)\to\mathbb EY^2=\sigma^2>0\) and
\(W_Y(x)\to\sigma^2\), because \(x^2q_Y(x)\leq\mathbb E(Y^2;|Y|>x)\to0\).
Hence \(H_Y(cx)/H_Y(x)\to\sigma^2/\sigma^2=1\) for every fixed
\(c>0\): \(V=H_Y\) is slowly varying, and
\cref{thm:p15-nda-normalizer-construction} supplies the scales
\(b_n\). Its defining equation reads \(b_n^2=nW_Y(b_n)\), and
\(b_n\to\infty\), so
\(b_n^2/n=W_Y(b_n)\to\sigma^2\), that is, \(b_n/\sqrt n\to\sigma\).
The same finite variance gives \(x^2q_X(x)/H_X(x)\to0\), since
\(x^2q_X(x)\leq\mathbb E(X^2;|X|>x)\to0\) while
\(H_X(x)\to\mathbb EX^2\geq\sigma^2>0\).
\Cref{thm:p15-nda-sufficiency} therefore applies, and the
centering it uses is the exact center \(n\mu\):
\[
 \frac{S_n-n\mu}{b_n}\Rightarrow\chi .
 \]
Keep the centering \(n\mu\) and replace the scale \(b_n\) by
\(\sqrt{n\mathbb EX^2}\). Then
\[
 \frac{S_n-n\mu}{\sqrt{n\mathbb EX^2}}
 =\frac{b_n}{\sqrt{n\mathbb EX^2}}\cdot\frac{S_n-n\mu}{b_n},
 \qquad
 \frac{b_n}{\sqrt{n\mathbb EX^2}}\to\frac{\sigma}{\sqrt{\mathbb EX^2}},
 \]
so by \cref{thm:p15-clt-perturbation}, with deterministic
\(A_n=b_n/\sqrt{n\mathbb EX^2}\) and \(B_n=0\), the left side converges weakly to
\((\sigma/\sqrt{\mathbb EX^2})\chi\), whose variance is
\(\sigma^2/\mathbb EX^2\). Since
\(\mathbb EX^2=\sigma^2+\mu^2\), this variance equals one exactly
when \(\mu=0\); for \(\mu\ne0\) it is strictly less than one, as stated.

For the logarithmic example, its density integrates to
\(2\int_1^\infty x^{-3}\,dx=1\), and its symmetric tails
give mean zero and infinite second moment.
Direct integration gives \(H_X(x)=2\log x\) and
\(q_X(x)=x^{-2}\) for \(x\geq1\), hence
\(W_X(x)=1+2\log x\). The root equation is
\[
 b_n^2=n(1+2\log b_n).
 \]
It implies \(b_n\geq\sqrt n\) eventually.
For each \(\delta>0\), \(1+2\log x\leq C_\delta x^\delta\)
eventually, by elementary integration or differentiation of
\((\log x)/x^\delta\). For \(0<\delta<2\) this gives
\(b_n\leq(C_\delta n)^{1/(2-\delta)}\) eventually.
Taking logarithms and then letting \(\delta\downarrow0\)
proves \(\log b_n/\log n\to1/2\).
Return to the exact equation:
\[
 b_n^2/(n\log n)
       =(1+2\log b_n)/\log n\to1.
 \]
The characterization and the resulting scale equivalence prove
the asserted \(\sqrt{n\log n}\) normalization.

For the power-tail family, let \(X\) be any symmetric finite
real random variable with \(\mathbb P(|X|>x)=x^{-\alpha}\) for
every \(x\geq1\), where \(0<\alpha<2\).

First, the hypothesis fixes the law of \(|X|\).
At \(x=1\) it gives \(\mathbb P(|X|>1)=1\), so
\(\mathbb P(|X|\leq1)=0\), and then
\(\mathbb P(|X|\leq x)\leq\mathbb P(|X|\leq1)=0\) for \(0\leq x<1\).
Thus the distribution function of \(|X|\) is
\[
 G(x)=\mathbb P(|X|\leq x)=
 \begin{cases}0,&x<1,\\ 1-x^{-\alpha},&x\geq1.\end{cases}
 \]
The law with density \(g(u)=\alpha u^{-1-\alpha}\mathbf1_{\{u>1\}}\)
has the same distribution function, because
\(\int_1^x\alpha u^{-1-\alpha}\,du=1-x^{-\alpha}\) for \(x\geq1\).
Two laws on \(\mathbb R\) with the same distribution function
coincide, so \(|X|\) has density \(g\).
Symmetry then fixes the law of \(X\) as well: for \(x\geq0\),
\(\mathbb P(X>x)=\mathbb P(X<-x)\) and these two events are
disjoint with union \(\{|X|>x\}\), so
\(\mathbb P(X>x)=\mathbb P(X<-x)=\tfrac12\mathbb P(|X|>x)\),
and \(\mathbb P(X=y)\leq\mathbb P(|X|=|y|)=0\) for every real
\(y\), because \(|X|\) has a density. Hence
\(\mathbb P(X\leq y)=1-\tfrac12\mathbb P(|X|>y)\) for \(y\geq0\)
and \(\mathbb P(X\leq y)=\mathbb P(X<y)=\tfrac12\mathbb P(|X|>-y)\)
for \(y<0\), which determines the distribution function of \(X\)
at every real point. Therefore the statement concerns exactly one law,
the symmetric law with density
\((\alpha/2)|x|^{-1-\alpha}\mathbf1_{\{|x|>1\}}\); that density is
symmetric, has total mass \(2\int_1^\infty(\alpha/2)u^{-1-\alpha}\,du=1\),
and has absolute tail \(x^{-\alpha}\) for \(x\geq1\), so the
law exists. The computations below use only the law of \(|X|\),
because \(X^2=|X|^2\) and \(q_X\) is a tail of \(|X|\).

Second, the truncated moments. For \(x\geq1\),
\(q_X(x)=x^{-\alpha}\) by hypothesis, and, since \(|X|>1\)
almost surely,
\[
 H_X(x)=\mathbb E(|X|^2;|X|\leq x)
       =\int_1^x u^2\,\alpha u^{-1-\alpha}\,du
       =\alpha\int_1^x u^{1-\alpha}\,du
       =\frac{\alpha}{2-\alpha}\bigl(x^{2-\alpha}-1\bigr),
 \]
where \(2-\alpha>0\) was used to integrate \(u^{1-\alpha}\).
This is strictly positive for \(x>1\).

Third, the ratio limit. For \(x>1\),
\[
 \frac{x^2q_X(x)}{H_X(x)}
 =\frac{x^{2-\alpha}}{\frac{\alpha}{2-\alpha}(x^{2-\alpha}-1)}
 =\frac{2-\alpha}{\alpha}\cdot\frac1{1-x^{-(2-\alpha)}} .
 \]
Because \(2-\alpha>0\), \(x^{-(2-\alpha)}\to0\) as
\(x\to\infty\), so the ratio tends to
\((2-\alpha)/\alpha\), which is strictly positive since
\(0<\alpha<2\).

Fourth, nondegeneracy, which
\cref{thm:p15-nda-characterization} requires. If \(X\) were
almost surely equal to a constant \(c\), symmetry would give
\(-X\) the same law, so \(c=-c\) and \(c=0\); but
\(\mathbb P(|X|>1)=1\). Hence the law is nondegenerate; indeed
\(\mathbb P(X>1)=\mathbb P(X<-1)=\tfrac12\).

Fifth, the conclusion. Suppose deterministic \(a_n\in\mathbb R\),
\(d_n>0\) gave \((S_n-a_n)/d_n\Rightarrow m+\tau\chi\) for some
\(m\in\mathbb R\) and \(\tau>0\), that is, a nondegenerate normal
limit. Then
\[
 \frac{S_n-(a_n+md_n)}{\tau d_n}
 =\frac1\tau\Bigl(\frac{S_n-a_n}{d_n}-m\Bigr)\Rightarrow\chi,
 \]
because the map \(y\mapsto(y-m)/\tau\) is continuous.
Condition 1 of \cref{thm:p15-nda-characterization} would hold
with centering \(a_n+md_n\) and scale \(\tau d_n\), and the
characterization, which applies because the law is nondegenerate,
would give condition 2, \(x^2q_X(x)/H_X(x)\to0\).
This contradicts the strictly positive limit
\((2-\alpha)/\alpha\). Hence no deterministic centering and
positive normalization yields a nondegenerate normal limit.

Finally, for \(X=1\), \(H_X(x)=1\) and \(q_X(x)=0\) at \(x>1\).
But every centered normalized sum is a constant, whose
characteristic function has modulus one; a nondegenerate
normal characteristic function has modulus less than one at
every nonzero argument. Such weak convergence is impossible.
For \(X=0\), both truncated second moment and tail vanish,
so the stated ratio is undefined. These examples explain the
nondegeneracy and denominator clauses directly.
\end{proof}
\end{annexitem}
